\section{Security, Privacy, and Evidence Integrity}

\subsection{Data boundary}

The characterized product posture is customer-hosted: decision data and generated row-level data
are intended to remain inside the deployment boundary, while review artifacts expose aggregates
and declared provenance. Actual network, identity, storage, backup, logging, and operator controls
remain deployment responsibilities and must be validated in the target environment.

This publication contains no customer data. The only row-level dataset in the companion is the
5,000-row generated balanced fixture used for utility reproduction. Its uncompressed SHA-256 is:
\begin{center}
  \IdentityText{b04f721d789226723066b3d6ae70e4ab2a3fa17c825ef1d0e04d7d779571982b}
\end{center}
The source manifest binds it to the exact v5.0.1 Gold artifact. It contains generated attributes,
not names, contact details, account identifiers, free text, credentials, or customer records.

\subsection{Bounded privacy observations}

The branch certificates report an exact-output wall and zero exact training-row reproductions in
the final amplification and intrinsic artifacts. That observation is bounded to the shared columns,
canonical predicate, and exact files in this run. It does not cover near matches, membership
inference, attribute inference, auxiliary information, linkage attacks, or future runs.

The reported privacy score is explicitly a quality heuristic. Near-duplicate status is
\texttt{not\_computed}; no differential-privacy mechanism, budget, accountant, or formal guarantee
is claimed. A deployment that requires a defined privacy risk threshold must add an approved
near-duplicate or distance-based protocol and, where appropriate, a separately designed formal
privacy mechanism.

\subsection{Four distinct integrity layers}

The evidence uses several mechanisms that should not be conflated:

\begin{enumerate}
  \item \textbf{File hashes} detect a byte change relative to a known digest.
  \item \textbf{Certificate content hashes and predecessor links} connect 21 bundled certificate
        files internally. Three legacy-name files repeat canonical content under the same content
        hash (their excluded signing timestamps differ), leaving 18 distinct hash nodes and 17
        distinct predecessor edges to one root. All 20 non-root file entries resolve, and the
        companion traverses every distinct node.
  \item \textbf{Workflow and artifact identities} bind the source repository, head commit, run,
        attempt, artifact ID, artifact API digest, inner bundle digest, and intake contract.
  \item \textbf{Candidate companion digest} binds all co-distributed evidence files to the value
        printed on the cover.
\end{enumerate}

The certificate JSON includes ECDSA algorithm, signature, timestamp, and key-fingerprint fields,
but the corresponding public key is not included in this companion. The offline verifier checks
their completeness and encoding, not their cryptographic authorship. The companion ZIP itself is
not signed. Therefore, an untrusted party could replace both a ZIP and an unauthenticated digest.
Independent verification requires the cover PDF or ZIP digest to arrive through a separately
trusted channel, or a future authorized release to add a verifiable external signature.

\subsection{Threats not resolved by this paper}

This characterization does not test deployment hardening, authorization policy, tenant isolation,
operator compromise, supply-chain compromise outside the recorded image, denial of service,
maliciously selected inputs, training-data poisoning, or exfiltration through a modified deployment.
It also does not prove that the public source and a privately operated image are identical beyond
the recorded digest relationship. Those controls belong in a deployment-specific security review.
